\documentclass[10pt]{article}  

%%%%%%%% PREÁMBULO %%%%%%%%%%%%
\title{Plantilla para prácticas de UGR}
\usepackage[spanish]{babel} %Indica que escribiermos en español
\usepackage[utf8]{inputenc} %Indica qué codificación se está usando ISO-8859-1(latin1)  o utf8  
\usepackage{amsmath} % Comandos extras para matemáticas (cajas para ecuaciones,
% etc)
\usepackage{amssymb} % Simbolos matematicos (por lo tanto)
\usepackage{graphicx} % Incluir imágenes en LaTeX
\usepackage{color} % Para colorear texto
\usepackage{subfigure} % subfiguras
\usepackage{float} %Podemos usar el especificador [H] en las figuras para que se
% queden donde queramos
\usepackage{capt-of} % Permite usar etiquetas fuera de elementos flotantes
% (etiquetas de figuras)
\usepackage{sidecap} % Para poner el texto de las imágenes al lado
	\sidecaptionvpos{figure}{c} % Para que el texto se alinie al centro vertical
\usepackage{caption} % Para poder quitar numeracion de figuras
\usepackage{commath} % funcionalidades extras para diferenciales, integrales,
% etc (\od, \dif, etc)
\usepackage{cancel} % para cancelar expresiones (\cancelto{0}{x})

\graphicspath{{/Users/jesusgarciamanday/Desktop/Master/CC-II/Practicas/Practica4/p4/Imagenes/}}

\usepackage{anysize} 					% Para personalizar el ancho de  los márgenes
\marginsize{2cm}{2cm}{2cm}{2cm} % Izquierda, derecha, arriba, abajo

\usepackage{appendix}
\renewcommand{\appendixname}{Apéndices}
\renewcommand{\appendixtocname}{Apéndices}
\renewcommand{\appendixpagename}{Apéndices} 

% Para que las referencias sean hipervínculos a las figuras o ecuaciones y
% aparezcan en color
\usepackage[colorlinks=true,plainpages=true,citecolor=blue,linkcolor=blue]{hyperref}
%\usepackage{hyperref} 
% Para agregar encabezado y pie de página
\usepackage{fancyhdr} 
\pagestyle{fancy}
\fancyhf{}
\fancyhead[L]{\footnotesize UGR} %encabezado izquierda
\fancyhead[R]{\footnotesize CCIA}   % dereecha
\fancyfoot[R]{\footnotesize Pr\'actica 4 - Hadoop }  % Pie derecha
\fancyfoot[C]{\thepage}  % centro
\fancyfoot[L]{\footnotesize Master en Ingenier\'ia Inform\'atica }  %izquierda
\renewcommand{\footrulewidth}{0.4pt}


\usepackage{listings} % Para usar código fuente
\definecolor{dkgreen}{rgb}{0,0.6,0} % Definimos colores para usar en el código
\definecolor{gray}{rgb}{0.5,0.5,0.5} 
% configuración para el lenguaje que queramos utilizar
\lstset{language=Matlab,
   keywords={break,case,catch,continue,else,elseif,end,for,function,
      global,if,otherwise,persistent,return,switch,try,while},
   basicstyle=\ttfamily,
   keywordstyle=\color{blue},
   commentstyle=\color{red},
   stringstyle=\color{dkgreen},
   numbers=left,
   numberstyle=\tiny\color{gray},
   stepnumber=1,
   numbersep=10pt,
   backgroundcolor=\color{white},
   tabsize=4,
   showspaces=false,
   showstringspaces=false}

\newcommand{\sen}{\operatorname{\sen}}	% Definimos el comando \sen para el seno
%en español

\title{Práctica 4 - Computación escalable y distribuida con Hadoop}

%%%%%%%% TERMINA PREÁMBULO %%%%%%%%%%%%

\begin{document}

%%%%%%%%%%%%%%%%%%%%%%%%%%%%%%%%%% PORTADA %%%%%%%%%%%%%%%%%%%%%%%%%%%%%%%%%%%%%%%%%%%%
																					%%%
\begin{center}																		%%%
\newcommand{\HRule}{\rule{\linewidth}{0.5mm}}									%%%\left
 																					%%%
\begin{minipage}{0.48\textwidth} \begin{flushleft}
%\includegraphics[scale = 0.63]{Imagenes/logo_upiita}
\end{flushleft}\end{minipage}
\begin{minipage}{0.48\textwidth} \begin{flushright}
%\includegraphics[scale = 0.35]{Imagenes/IPN}
\end{flushright}\end{minipage}

													 								%%%
\vspace*{-1.5cm}								%%%
																					%%%	
\textsc{\huge Universidad de\\ \vspace{5px} Granada}\\[1.5cm]	

\textsc{\LARGE Master Profesional en Ingenier\'ia Inform\'atica }\\[1.5cm]													%%%

\begin{minipage}{0.9\textwidth} 
\begin{center}																					%%%
\textsc{\LARGE Pr\'actica 4}
\end{center}
\end{minipage}\\[0.5cm]
%%%
    																				%%%
 			\vspace*{1cm}																		%%%
																					%%%
\HRule \\[0.4cm]																	%%%
{ \huge \bfseries Hadoop}\\[0.4cm]	%%%
 																					%%%
\HRule \\[1.5cm]																	%%%
 																				%%%
																					%%%
\begin{minipage}{0.46\textwidth}													%%%
\begin{flushleft} \large															%%%
\emph{Autor:}\\	
Manuel Jes\'us Garc\'ia Manday (nickter@correo.ugr.es)\\
%%%
			%\vspace*{2cm}	
            													%%%
										 						%%%
\end{flushleft}																		%%%
\end{minipage}		
																%%%
\begin{minipage}{0.52\textwidth}		
\vspace{-0.6cm}											%%%
\begin{flushright} \large															%%%
													%%%
\end{flushright}																	%%%
\end{minipage}	
\vspace*{1cm}
%\begin{flushleft}
 	
%\end{flushleft}
%%%
 		\flushleft{\textbf{\Large Master en Ingenier\'ia Inform\'atica}	}\\																		%%%
\vspace{2cm} 																				
\begin{center}																					
{\large \today}																	%%%
 			\end{center}												  						
\end{center}							 											
																					
\newpage																		
%%%%%%%%%%%%%%%%%%%% TERMINA PORTADA %%%%%%%%%%%%%%%%%%%%%%%%%%%%%%%%

\tableofcontents 

\newpage

\section{Objetivo.}
El objetivo de esta práctica es realizar programas escalables para mejorar la eficiencia en entornos Big Data.\\


\section{Introduccíon.}
Para comenzar a realizar las tareas  que se piden en esta práctica, es necesario en primer lugar realizar una serie de pasos iniciales que se describen a continuación. \\

Realizamos una conexión remota hacía el servidor \textbf{haddop.ugr.es} y una vez dentro creamos una carpeta nueva donde descargaremos el código Java de los programas. Comprobamos tambien que los datos de entrada se encuentran disponibles.\\

 \begin{figure}[H]
	\begin{center}
 		\includegraphics[width = 1.00\textwidth]{p4-img1}
 		\captionof{figure}{\label{fig:IPN}Conexion remota a \textbf{hadoop.ugr.es}.} 
	\end{center} 
\end{figure}

 \begin{figure}[H]
	\begin{center}
 		\includegraphics[width = 1.00\textwidth]{p4-img2}
 		\captionof{figure}{\label{fig:IPN}Copiando los ficheros Java.} 
	\end{center} 
\end{figure}

 \begin{figure}[H]
	\begin{center}
 		\includegraphics[width = 1.00\textwidth]{p4-img3}
 		\captionof{figure}{\label{fig:IPN}Datos de entrada.} 
	\end{center} 
\end{figure}

Viendo que están disponibles, ahora nos creamos un directorio local para las clases de java. \\

 \begin{figure}[H]
	\begin{center}
 		\includegraphics[width = 1.00\textwidth]{p4-img4}
 		\captionof{figure}{\label{fig:IPN}Directorio local.} 
	\end{center} 
\end{figure}

Con la preconfiguración realizada pasamos a realizar las diferentes tareas que es exponen en la práctica. \\


\section{Calcula el valor mínimo de la variable (columna) 5.} 
La primera de ellas es calcular el mínimo sobre el conjunto de valores del dataset, por lo que una vez que tenemos todos los ficheros java correspondientes, ahora toca compilarlos para crear el fichero \textbf{.jar} a continuación y ejecutarlo en \textbf{hadoop}.\\

\begin{figure}[H]
	\begin{center}
 		\includegraphics[width = 1.00\textwidth]{p4-img5}
 		\captionof{figure}{\label{fig:IPN}Compilamos y ejecutamos (I).} 
	\end{center} 
\end{figure}

\begin{figure}[H]
	\begin{center}
 		\includegraphics[width = 1.00\textwidth]{p4-img6}
 		\captionof{figure}{\label{fig:IPN}Compilamos y ejecutamos (II).} 
	\end{center} 
\end{figure}

\begin{figure}[H]
	\begin{center}
 		\includegraphics[width = 1.00\textwidth]{p4-img7}
 		\captionof{figure}{\label{fig:IPN}Compilamos y ejecutamos (III).} 
	\end{center} 
\end{figure}

\begin{figure}[H]
	\begin{center}
 		\includegraphics[width = 1.00\textwidth]{p4-img8}
 		\captionof{figure}{\label{fig:IPN}Compilamos y ejecutamos (IV).} 
	\end{center} 
\end{figure}

Por último comprobamos el resultado para ver si se ha realizado correctamente. \\

\begin{figure}[H]
	\begin{center}
 		\includegraphics[width = 1.00\textwidth]{p4-img9}
 		\captionof{figure}{\label{fig:IPN}Comprobando resultado.} 
	\end{center} 
\end{figure}


\section{Calcula el valor máximo de la variable (columna) 5.} 
Para calcular el valor máximo de la variable (columna) 5 vamos a crear clases de java correspondientes para hacer dicha operación. Comenzaremos por crear la clase \textbf{MaxMapper} con su fichero java correspondiente. \\

\begin{figure}[H]
	\begin{center}
 		\includegraphics[width = 1.00\textwidth]{p4-img10}
 		\captionof{figure}{\label{fig:IPN}Clase MaxMapper.} 
	\end{center} 
\end{figure}

A continuación creamos la clase correspondientes para el Reducer (\textbf{MaxReducer}) y la clase principal \textbf{Max} donde estará el main. \\

\begin{figure}[H]
	\begin{center}
 		\includegraphics[width = 1.00\textwidth]{p4-img11}
 		\captionof{figure}{\label{fig:IPN}Clase MaxReducer.} 
	\end{center} 
\end{figure}

\begin{figure}[H]
	\begin{center}
 		\includegraphics[width = 1.00\textwidth]{p4-img12}
 		\captionof{figure}{\label{fig:IPN}Clase Max.} 
	\end{center} 
\end{figure}

Una vez que tenemos todos los ficheros procedemos a realizar los mismos pasos que con el ejercicio anterior compilando las clases, creando el fichero jar y ejecutandolo en \textbf{hadoop}. \\

\begin{figure}[H]
	\begin{center}
 		\includegraphics[width = 1.00\textwidth]{p4-img13}
 		\captionof{figure}{\label{fig:IPN}Compilamos y ejecutamos (I).} 
	\end{center} 
\end{figure}

\begin{figure}[H]
	\begin{center}
 		\includegraphics[width = 1.00\textwidth]{p4-img14}
 		\captionof{figure}{\label{fig:IPN}Compilamos y ejecutamos (II).}  
	\end{center} 
\end{figure}

\begin{figure}[H]
	\begin{center}
 		\includegraphics[width = 1.00\textwidth]{p4-img15}
 		\captionof{figure}{\label{fig:IPN}Compilamos y ejecutamos (III).} 
	\end{center} 
\end{figure}

\begin{figure}[H]
	\begin{center}
 		\includegraphics[width = 1.00\textwidth]{p4-img16}
 		\captionof{figure}{\label{fig:IPN}Compilamos y ejecutamos (IV).} 
	\end{center} 
\end{figure}

Comprobamos que el resultado nos arroja el valor máximo. \\

\begin{figure}[H]
	\begin{center}
 		\includegraphics[width = 1.00\textwidth]{p4-img17}
 		\captionof{figure}{\label{fig:IPN}Comprobando resultado.} 
	\end{center} 
\end{figure}


\section{Calcula al mismo tiempo los valores máximo y mínimo de la variable 5.} 
Al igual que para los anteriores cálculos creamos los ficheros correspondientes, aunque el fichero con la función \textbf{Mapper} y el principal con el main solo cambia el nombre, en la clase con la función \textbf{Reducer} es donde se realizan los cambios necesarios para obtener los resultados esperados. \\

\begin{figure}[H]
	\begin{center}
 		\includegraphics[width = 1.00\textwidth]{p4-img19}
 		\captionof{figure}{\label{fig:IPN}Clase MaxMinMapper.} 
	\end{center} 
\end{figure}

\begin{figure}[H]
	\begin{center}
 		\includegraphics[width = 1.00\textwidth]{p4-img18}
 		\captionof{figure}{\label{fig:IPN}Clase MaxMin.} 
	\end{center} 
\end{figure}

\begin{figure}[H]
	\begin{center}
 		\includegraphics[width = 1.00\textwidth]{p4-img21}
 		\captionof{figure}{\label{fig:IPN}Clase MaxMinReducer.} 
	\end{center} 
\end{figure}

Ahora toca compilarlos para adjuntarlos en el fichero \textbf{jar} y realizar la ejecución en \textbf{hadoop}. \\

\begin{figure}[H]
	\begin{center}
 		\includegraphics[width = 1.00\textwidth]{p4-img20}
 		\captionof{figure}{\label{fig:IPN}Compilamos y ejecutamos (I).} 
	\end{center} 
\end{figure}

\begin{figure}[H]
	\begin{center}
 		\includegraphics[width = 1.00\textwidth]{p4-img22}
 		\captionof{figure}{\label{fig:IPN}Compilamos y ejecutamos (II).}  
	\end{center} 
\end{figure}

\begin{figure}[H]
	\begin{center}
 		\includegraphics[width = 1.00\textwidth]{p4-img23}
 		\captionof{figure}{\label{fig:IPN}Compilamos y ejecutamos (III).} 
	\end{center} 
\end{figure}

\begin{figure}[H]
	\begin{center}
 		\includegraphics[width = 1.00\textwidth]{p4-img24}
 		\captionof{figure}{\label{fig:IPN}Compilamos y ejecutamos (IV).} 
	\end{center} 
\end{figure}


Por último comprobamos el resultado obtenido. \\

\begin{figure}[H]
	\begin{center}
 		\includegraphics[width = 1.00\textwidth]{p4-img25}
 		\captionof{figure}{\label{fig:IPN}Comprobando resultado.} 
	\end{center} 
\end{figure}


\section{Calcula los valores máximo y mínimo de todas las variables (salvo la última, que es la  etiqueta de la clase).} 
Para este cálculo solo necesitaremos adaptar la clase con la función \textbf{Mapper} para cada columna ya que el \textbf{Reducer} es el mismo. Ahora la función \textbf{Mapper} asignará diferentes key, uno por cada columna. \\

\begin{figure}[H]
	\begin{center}
 		\includegraphics[width = 1.00\textwidth]{p4-img26}
 		\captionof{figure}{\label{fig:IPN}Clase MaxMinMapper.} 
	\end{center} 
\end{figure}

Como en los anteriores ejercicios, compilamos y ejecutamos las clases. \\

\begin{figure}[H]
	\begin{center}
 		\includegraphics[width = 1.00\textwidth]{p4-img27}
 		\captionof{figure}{\label{fig:IPN}Compilamos y ejecutamos (I).} 
	\end{center} 
\end{figure}

\begin{figure}[H]
	\begin{center}
 		\includegraphics[width = 1.00\textwidth]{p4-img28}
 		\captionof{figure}{\label{fig:IPN}Compilamos y ejecutamos (II).} 
	\end{center} 
\end{figure}

\begin{figure}[H]
	\begin{center}
 		\includegraphics[width = 1.00\textwidth]{p4-img29}
 		\captionof{figure}{\label{fig:IPN}Compilamos y ejecutamos (III).} 
	\end{center} 
\end{figure}

Vemos el resultado que se ha obtenido. \\

\begin{figure}[H]
	\begin{center}
 		\includegraphics[width = 1.00\textwidth]{p4-img30}
 		\captionof{figure}{\label{fig:IPN}Comprobando resultado.} 
	\end{center} 
\end{figure}


\section{Realizar la media de la variable 5.} 
Para realizar esta tarea será necesario crear dos funciones nuevas para \textbf{Mapper} y \textbf{Reducer} con sus respectivas clases ya que difieren de las tareas anteriores. \\

\begin{figure}[H]
	\begin{center}
 		\includegraphics[width = 1.00\textwidth]{p4-img31}
 		\captionof{figure}{\label{fig:IPN}Clase Average.} 
	\end{center} 
\end{figure}

\begin{figure}[H]
	\begin{center}
 		\includegraphics[width = 1.00\textwidth]{p4-img32}
 		\captionof{figure}{\label{fig:IPN}Clase AverageMapper.} 
	\end{center} 
\end{figure}

\begin{figure}[H]
	\begin{center}
 		\includegraphics[width = 1.00\textwidth]{p4-img33}
 		\captionof{figure}{\label{fig:IPN}Clase AverageReducer.} 
	\end{center} 
\end{figure}

Con las funciones creadas en su correspondiente clase pasamos a compilar y a ejecutar. \\

\begin{figure}[H]
	\begin{center}
 		\includegraphics[width = 1.00\textwidth]{p4-img34}
 		\captionof{figure}{\label{fig:IPN}Compilamos y ejecutamos (I).} 
	\end{center} 
\end{figure}

\begin{figure}[H]
	\begin{center}
 		\includegraphics[width = 1.00\textwidth]{p4-img35}
 		\captionof{figure}{\label{fig:IPN}Compilamos y ejecutamos ()I).} 
	\end{center} 
\end{figure}

\begin{figure}[H]
	\begin{center}
 		\includegraphics[width = 1.00\textwidth]{p4-img36}
 		\captionof{figure}{\label{fig:IPN}Compilamos y ejecutamos (III).} 
	\end{center} 
\end{figure}

Y por último comprobamos el resultado. \\

\begin{figure}[H]
	\begin{center}
 		\includegraphics[width = 1.00\textwidth]{p4-img37}
 		\captionof{figure}{\label{fig:IPN}Comprobando resultado.} 
	\end{center} 
\end{figure}


\section{Obtener la media de todas las variables (salvo la clase).} 
En esta tarea solo necesitamos modificar la función \textbf{Mapper} con respecto al ejercicio anterior donde la función \textbf{Reduce} será la misma, por lo que será necesario modificar únicamente un fichero. \\

\begin{figure}[H]
	\begin{center}
 		\includegraphics[width = 1.00\textwidth]{p4-img38}
 		\captionof{figure}{\label{fig:IPN}Clase AverageMapper.} 
	\end{center} 
\end{figure}

Una vez modificado el método \textbf{Mapper} podemos compilar y ejecutar. \\

\begin{figure}[H]
	\begin{center}
 		\includegraphics[width = 1.00\textwidth]{p4-img39}
 		\captionof{figure}{\label{fig:IPN}Compilamos y ejecutamos (I).} 
	\end{center} 
\end{figure}

\begin{figure}[H]
	\begin{center}
 		\includegraphics[width = 1.00\textwidth]{p4-img40}
 		\captionof{figure}{\label{fig:IPN}Compilamos y ejecutamos (II).} 
	\end{center} 
\end{figure}

Sólo nos queda comprobar el resultado. \\

\begin{figure}[H]
	\begin{center}
 		\includegraphics[width = 1.00\textwidth]{p4-img41}
 		\captionof{figure}{\label{fig:IPN}Comprobando resultado.} 
	\end{center} 
\end{figure}


\section{Comprobar si el conjunto de datos ECBDL es balanceado o no balanceado, es decir, que el ratio entre clases sea menor o mayor que 1.5 respectivamente.} 
Para calcular si el conjunto de datos es balanceado o no balanceado a partir del ratio entre las clases es necesario definirse dos nuevas funciones \textbf{Mapper} y \textbf{Reducer} como se muestra en las imágenes. \\

\begin{figure}[H]
	\begin{center}
 		\includegraphics[width = 1.00\textwidth]{p4-img42}
 		\captionof{figure}{\label{fig:IPN}Clase Balancer.} 
	\end{center} 
\end{figure}

\begin{figure}[H]
	\begin{center}
 		\includegraphics[width = 1.00\textwidth]{p4-img43}
 		\captionof{figure}{\label{fig:IPN}Clase BalancerMapper.} 
	\end{center} 
\end{figure}

\begin{figure}[H]
	\begin{center}
 		\includegraphics[width = 1.00\textwidth]{p4-img44}
 		\captionof{figure}{\label{fig:IPN}ClaseBalancerReducer.} 
	\end{center} 
\end{figure}

De nuevo como en las tareas anteriores compilamos y ejecutamos. \\

\begin{figure}[H]
	\begin{center}
 		\includegraphics[width = 1.00\textwidth]{p4-img45}
 		\captionof{figure}{\label{fig:IPN}Compilamos y ejecutamos (I).} 
	\end{center} 
\end{figure}

\begin{figure}[H]
	\begin{center}
 		\includegraphics[width = 1.00\textwidth]{p4-img46}
 		\captionof{figure}{\label{fig:IPN}Compilamos y ejecutamos (II).} 
	\end{center} 
\end{figure}

\begin{figure}[H]
	\begin{center}
 		\includegraphics[width = 1.00\textwidth]{p4-img47}
 		\captionof{figure}{\label{fig:IPN}Compilamos y ejecutamos (III).} 
	\end{center} 
\end{figure}

Y comprobamos el resultado donde se puede ver que el ratio es mucho mayor que 1.5 por lo que las clases no están balanceadas. \\

\begin{figure}[H]
	\begin{center}
 		\includegraphics[width = 1.00\textwidth]{p4-img48}
 		\captionof{figure}{\label{fig:IPN}Comprobando resultado.} 
	\end{center} 
\end{figure}


\section{Cálculo del coeficiente de correlación entre todas las parejas de variables.} 
Creamos los ficheros correspondientes donde se definen las funciones \textbf{Mapper} y \textbf{Reducer} para calcular la correlación entre cada par de variables del conjunto de datos. \\

\begin{figure}[H]
	\begin{center}
 		\includegraphics[width = 1.00\textwidth]{p4-img49}
 		\captionof{figure}{\label{fig:IPN}Clase Correlacion.} 
	\end{center} 
\end{figure}

\begin{figure}[H]
	\begin{center}
 		\includegraphics[width = 1.00\textwidth]{p4-img50}
 		\captionof{figure}{\label{fig:IPN}Clase CorrelacionMapper.} 
	\end{center} 
\end{figure}

\begin{figure}[H]
	\begin{center}
 		\includegraphics[width = 1.00\textwidth]{p4-img51}
 		\captionof{figure}{\label{fig:IPN}Clase CorrelacionReducer.} 
	\end{center} 
\end{figure}

Una vez que tenemos los ficheros con las funciones implementadas, pasamos a compilar las clases para empaquetarlas en el archivo \textbf{jar} correspondiente y ejecutarlo en \textbf{hadoop}.\\

\begin{figure}[H]
	\begin{center}
 		\includegraphics[width = 1.00\textwidth]{p4-img52}
 		\captionof{figure}{\label{fig:IPN}Compilamos y ejecutamos (I).} 
	\end{center} 
\end{figure}

\begin{figure}[H]
	\begin{center}
 		\includegraphics[width = 1.00\textwidth]{p4-img53}
 		\captionof{figure}{\label{fig:IPN}Compilamos y ejecutamos (II).} 
	\end{center} 
\end{figure}

\begin{figure}[H]
	\begin{center}
 		\includegraphics[width = 1.00\textwidth]{p4-img54}
 		\captionof{figure}{\label{fig:IPN}Compilamos y ejecutamos (III).} 
	\end{center} 
\end{figure}

\begin{figure}[H]
	\begin{center}
 		\includegraphics[width = 1.00\textwidth]{p4-img55}
 		\captionof{figure}{\label{fig:IPN}Compilamos y ejecutamos (IV).} 
	\end{center} 
\end{figure}

\begin{figure}[H]
	\begin{center}
 		\includegraphics[width = 1.00\textwidth]{p4-img56}
 		\captionof{figure}{\label{fig:IPN}Compilamos y ejecutamos (V).} 
	\end{center} 
\end{figure}

Con la ejecución terminada correctamente pasamos a ver los resultados obtenidos. \\

\begin{figure}[H]
	\begin{center}
 		\includegraphics[width = 1.00\textwidth]{p4-img57}
 		\captionof{figure}{\label{fig:IPN}Comprobando resultado.} 
	\end{center} 
\end{figure}


\section{Tareas extras}
\subsection{Parametrizar la columna sobre la que se quiere realizar el cálculo estadístico.}
En esta tarea se parametrizará la columna sobre la que se va a realizar el cálculo del valor mínimo, permitiendo introducir desde consola el número de la consola. Las clase principal y la clase \textbf{Mapper} es donde se han reflejado los cambios, ya que se ha utilizado el objeto \textbf{Configuration} para pasar el valor de la columna como parámetro. \\

\begin{figure}[H]
	\begin{center}
 		\includegraphics[width = 1.00\textwidth]{p4-img67}
 		\captionof{figure}{\label{fig:IPN}Clase MinParam.} 
	\end{center} 
\end{figure}

\begin{figure}[H]
	\begin{center}
 		\includegraphics[width = 1.00\textwidth]{p4-img68}
 		\captionof{figure}{\label{fig:IPN}Clase MinParamMapper.} 
	\end{center} 
\end{figure}

\begin{figure}[H]
	\begin{center}
 		\includegraphics[width = 1.00\textwidth]{p4-img69}
 		\captionof{figure}{\label{fig:IPN}Clase MinParamReducer.} 
	\end{center} 
\end{figure}

Compilamos las ficheros y ejecutamos pasando por consola la columna 5. \\

\begin{figure}[H]
	\begin{center}
 		\includegraphics[width = 1.00\textwidth]{p4-img70}
 		\captionof{figure}{\label{fig:IPN}Compilamos y ejecutamos con la columna 5 (I).} 
	\end{center} 
\end{figure}

\begin{figure}[H]
	\begin{center}
 		\includegraphics[width = 1.00\textwidth]{p4-img71}
 		\captionof{figure}{\label{fig:IPN}Compilamos y ejecutamos con la columna 5 (II).} 
	\end{center} 
\end{figure}

\begin{figure}[H]
	\begin{center}
 		\includegraphics[width = 1.00\textwidth]{p4-img72}
 		\captionof{figure}{\label{fig:IPN}Compilamos y ejecutamos con la columna 5 (III).} 
	\end{center} 
\end{figure}

\begin{figure}[H]
	\begin{center}
 		\includegraphics[width = 1.00\textwidth]{p4-img73}
 		\captionof{figure}{\label{fig:IPN}Compilamos y ejecutamos con la columna 5 (IV).} 
	\end{center} 
\end{figure}

Ahora lo ejecutamos con la columna 1. \\

\begin{figure}[H]
	\begin{center}
 		\includegraphics[width = 1.00\textwidth]{p4-img74}
 		\captionof{figure}{\label{fig:IPN}Compilamos y ejecutamos con la columna 1 (I).} 
	\end{center} 
\end{figure}

\begin{figure}[H]
	\begin{center}
 		\includegraphics[width = 1.00\textwidth]{p4-img75}
 		\captionof{figure}{\label{fig:IPN}Compilamos y ejecutamos con la columna 1 (II).} 
	\end{center} 
\end{figure}

\begin{figure}[H]
	\begin{center}
 		\includegraphics[width = 1.00\textwidth]{p4-img76}
 		\captionof{figure}{\label{fig:IPN}Compilamos y ejecutamos con la columna 1 (III).} 
	\end{center} 
\end{figure}

\subsection{Combinar el cálculo de todos los estadístico en una función.}
Para esta tarea se van a combinar todos los cálculos estadísticos sobre la función \textbf{Reducer} para una variable. \\

\begin{figure}[H]
	\begin{center}
 		\includegraphics[width = 1.00\textwidth]{p4-img88}
 		\captionof{figure}{\label{fig:IPN}Clase CombStat.} 
	\end{center} 
\end{figure}

\begin{figure}[H]
	\begin{center}
 		\includegraphics[width = 1.00\textwidth]{p4-img89}
 		\captionof{figure}{\label{fig:IPN}Clase CombStatMapper.} 
	\end{center} 
\end{figure}

\begin{figure}[H]
	\begin{center}
 		\includegraphics[width = 1.00\textwidth]{p4-img90}
 		\captionof{figure}{\label{fig:IPN}Clase CombStatReducer (I).} 
	\end{center} 
\end{figure}

\begin{figure}[H]
	\begin{center}
 		\includegraphics[width = 1.00\textwidth]{p4-img91}
 		\captionof{figure}{\label{fig:IPN}Clase CombStatReducer (II).} 
	\end{center} 
\end{figure}





\subsection{Calcular los datos estadísticos sobre todas las columnas.}
En este ejercicio vamos a calcular la media, la covarianza, la desviación típica y el coeficiente de correlación sobre todas las variables.

\begin{figure}[H]
	\begin{center}
 		\includegraphics[width = 1.00\textwidth]{p4-img77}
 		\captionof{figure}{\label{fig:IPN}Clase Stat.} 
	\end{center} 
\end{figure}

\begin{figure}[H]
	\begin{center}
 		\includegraphics[width = 1.00\textwidth]{p4-img78}
 		\captionof{figure}{\label{fig:IPN}Clase StatMapper.} 
	\end{center} 
\end{figure}

\begin{figure}[H]
	\begin{center}
 		\includegraphics[width = 1.00\textwidth]{p4-img79}
 		\captionof{figure}{\label{fig:IPN}Clase StatReducer (I).} 
	\end{center} 
\end{figure}

\begin{figure}[H]
	\begin{center}
 		\includegraphics[width = 1.00\textwidth]{p4-img80}
 		\captionof{figure}{\label{fig:IPN}Clase StatReducer (II).} 
	\end{center} 
\end{figure}

Compilamos y ejecutamos para ver el tiempo empleado en realizar los cálculos. \\

\begin{figure}[H]
	\begin{center}
 		\includegraphics[width = 1.00\textwidth]{p4-img81}
 		\captionof{figure}{\label{fig:IPN}Compilamos y ejecutamos (I).} 
	\end{center} 
\end{figure}

\begin{figure}[H]
	\begin{center}
 		\includegraphics[width = 1.00\textwidth]{p4-img93}
 		\captionof{figure}{\label{fig:IPN}Compilamos y ejecutamos (II).} 
	\end{center} 
\end{figure}

\begin{figure}[H]
	\begin{center}
 		\includegraphics[width = 1.00\textwidth]{p4-img83}
 		\captionof{figure}{\label{fig:IPN}Compilamos y ejecutamos (III).} 
	\end{center} 
\end{figure}

\begin{figure}[H]
	\begin{center}
 		\includegraphics[width = 1.00\textwidth]{p4-img84}
 		\captionof{figure}{\label{fig:IPN}Compilamos y ejecutamos (IV).} 
	\end{center} 
\end{figure}

\begin{figure}[H]
	\begin{center}
 		\includegraphics[width = 1.00\textwidth]{p4-img85}
 		\captionof{figure}{\label{fig:IPN}Compilamos y ejecutamos (V).} 
	\end{center} 
\end{figure}

\begin{figure}[H]
	\begin{center}
 		\includegraphics[width = 1.00\textwidth]{p4-img86}
 		\captionof{figure}{\label{fig:IPN}Compilamos y ejecutamos (VI).} 
	\end{center} 
\end{figure}

\begin{figure}[H]
	\begin{center}
 		\includegraphics[width = 1.00\textwidth]{p4-img92}
 		\captionof{figure}{\label{fig:IPN}Compilamos y ejecutamos (VII).} 
	\end{center} 
\end{figure}

En esta última imagen vemos el tiempo empleado con este conjunto de datos que lo compararemos en el ejercicio siguiente.\\


\subsection{Calcular los datos estadísticos sobre todas las columnas sobre un conjunto de mayor volumen.}
En este ejercicio realizaremos las mismas pruebas que en el anterior pero sobre un conjunto de datos mayor para comprobar la diferencia del tiempo empleado. \\

\begin{figure}[H]
	\begin{center}
 		\includegraphics[width = 1.00\textwidth]{p4-img77}
 		\captionof{figure}{\label{fig:IPN}Clase Stat.} 
	\end{center} 
\end{figure}

\begin{figure}[H]
	\begin{center}
 		\includegraphics[width = 1.00\textwidth]{p4-img78}
 		\captionof{figure}{\label{fig:IPN}Clase StatMapper.} 
	\end{center} 
\end{figure}

\begin{figure}[H]
	\begin{center}
 		\includegraphics[width = 1.00\textwidth]{p4-img79}
 		\captionof{figure}{\label{fig:IPN}Clase StatReducer (I).} 
	\end{center} 
\end{figure}

\begin{figure}[H]
	\begin{center}
 		\includegraphics[width = 1.00\textwidth]{p4-img80}
 		\captionof{figure}{\label{fig:IPN}Clase StatReducer (II).} 
	\end{center} 
\end{figure}

Compilamos y ejecutamos para ver el tiempo empleado en realizar los cálculos. \\

\begin{figure}[H]
	\begin{center}
 		\includegraphics[width = 1.00\textwidth]{p4-img81}
 		\captionof{figure}{\label{fig:IPN}Compilamos y ejecutamos (I).} 
	\end{center} 
\end{figure}

\begin{figure}[H]
	\begin{center}
 		\includegraphics[width = 1.00\textwidth]{p4-img82}
 		\captionof{figure}{\label{fig:IPN}Compilamos y ejecutamos (II).} 
	\end{center} 
\end{figure}

\begin{figure}[H]
	\begin{center}
 		\includegraphics[width = 1.00\textwidth]{p4-img83}
 		\captionof{figure}{\label{fig:IPN}Compilamos y ejecutamos (III).} 
	\end{center} 
\end{figure}

\begin{figure}[H]
	\begin{center}
 		\includegraphics[width = 1.00\textwidth]{p4-img84}
 		\captionof{figure}{\label{fig:IPN}Compilamos y ejecutamos (IV).} 
	\end{center} 
\end{figure}

\begin{figure}[H]
	\begin{center}
 		\includegraphics[width = 1.00\textwidth]{p4-img85}
 		\captionof{figure}{\label{fig:IPN}Compilamos y ejecutamos (V).} 
	\end{center} 
\end{figure}

\begin{figure}[H]
	\begin{center}
 		\includegraphics[width = 1.00\textwidth]{p4-img86}
 		\captionof{figure}{\label{fig:IPN}Compilamos y ejecutamos (VI).} 
	\end{center} 
\end{figure}

\begin{figure}[H]
	\begin{center}
 		\includegraphics[width = 1.00\textwidth]{p4-img87}
 		\captionof{figure}{\label{fig:IPN}Compilamos y ejecutamos (VII).} 
	\end{center} 
\end{figure}

Se puede ver como el tiempo ha aumentado hasta tres veces sobre el conjunto anterior como se comprueba con las imagenes.\\

\subsection{Acelerar el proceso de cómputo descargando al Reducer de parte de la tarea.}
En este ejercicio se calcula de nuevo el mínimo sobre la variable 5 pero utilizando el proceso \textbf{Cleanup} para descargar al \textbf{Reducer} de parte de la tarea. En la siguientes imágenes se muestran las clases con las funciones necesarias.\\

\begin{figure}[H]
	\begin{center}
 		\includegraphics[width = 1.00\textwidth]{p4-img58}
 		\captionof{figure}{\label{fig:IPN}Clase MinOpt.} 
	\end{center} 
\end{figure}

\begin{figure}[H]
	\begin{center}
 		\includegraphics[width = 1.00\textwidth]{p4-img59}
 		\captionof{figure}{\label{fig:IPN}Clase MinOptMapper.} 
	\end{center} 
\end{figure}

\begin{figure}[H]
	\begin{center}
 		\includegraphics[width = 1.00\textwidth]{p4-img60}
 		\captionof{figure}{\label{fig:IPN}Clase MinOptReducer.} 
	\end{center} 
\end{figure}

Compilamos y ejecutamos una vez que tenemos definidas las funciones \textbf{Mapper}, \textbf{Cleanup} y \textbf{Reducer}. \\

\begin{figure}[H]
	\begin{center}
 		\includegraphics[width = 1.00\textwidth]{p4-img61}
 		\captionof{figure}{\label{fig:IPN}Compilamos y ejecutamos (I).} 
	\end{center} 
\end{figure}

\begin{figure}[H]
	\begin{center}
 		\includegraphics[width = 1.00\textwidth]{p4-img62}
 		\captionof{figure}{\label{fig:IPN}Compilamos y ejecutamos (II).} 
	\end{center} 
\end{figure}

\begin{figure}[H]
	\begin{center}
 		\includegraphics[width = 1.00\textwidth]{p4-img63}
 		\captionof{figure}{\label{fig:IPN}Compilamos y ejecutamos (III).} 
	\end{center} 
\end{figure}

Y por último comprobamos el resultado donde podemos ver que se obtiene el mismo que en el ejercicio anterior pero de esta manera el proceso \textbf{MapReduce} ha sido más rápido debido a que mucha de la tarea que realizaba en el ejercicio anterior la función \textbf{Reducer} en este se ha trasladado al \textbf{Mapper}. La diferencia de tiempo empleado en la CPU es considerable como se puede ver en las imágenes anteriores. \\

 \begin{figure}[H]
	\begin{center}
 		\includegraphics[width = 1.00\textwidth]{p4-img64}
 		\captionof{figure}{\label{fig:IPN}Comprobando resultado.} 
	\end{center} 
\end{figure}

\begin{figure}[H]
	\begin{center}
 		\includegraphics[width = 1.00\textwidth]{p4-img65}
 		\captionof{figure}{\label{fig:IPN}MapReduce sin proceso Cleanup.} 
	\end{center} 
\end{figure}

\begin{figure}[H]
	\begin{center}
 		\includegraphics[width = 1.00\textwidth]{p4-img66}
 		\captionof{figure}{\label{fig:IPN}Compilamos con proceso Cleanup.} 
	\end{center} 
\end{figure}

\end{document}